% ============================================================================
% Chapter 24 --- Wave 6: Endpoint and posture
% ----------------------------------------------------------------------------
% Purpose : extend TAC enforcement to managed and BYOD endpoints. Must
%           surface the OSS gap on iOS/Android (D11.3) as a known
%           limitation.
% Page target: ~5 pages.
% Dependencies: Waves 1, 4 required; Waves 2, 5 recommended (D17.3);
%               reference layer Ch. 11 (unified posture schema D11.6,
%               BYOD via attachment D11.7).
% Mirrors wave template (D18.2).
% ----------------------------------------------------------------------------
% Decisions registered in _coherence-EN.md:
%   D24.1 posture signals advisory before enforcing (parallel to Wave 4 shadow)
%   D24.2 unified posture schema is the wave's central artefact
%   D24.3 OSS gap on iOS/Android surfaces here as explicit limitation
%   D24.4 BYOD posture remains network-side per D11.7
%   D24.5 exit criteria reference Endpoint conformance subset (D16.5)
% ============================================================================

\chapter{Wave~6 --- Endpoint and posture}
\label{ch:wave6-endpoint}

\begin{caveatbox}[Dependencies]
\noindent Wave~6 applies the wave chapter template established in
\trchap{ch:wave0-foundations}{}. Required pre-requisites are
Wave~1 (identity provider emitting the institutional schema) and
Wave~4 (policy plane consuming attribute and posture signals);
Waves~2 (observability) and 5 (pilot) strengthen the wave when
present. It instantiates the endpoint contract of
\trchap{ch:endpoint}{}, including the unified posture schema of
D11.6 and the boundary with network attachment of D11.7. Closure
is gated by the three Endpoint tests of \trchap{ch:conformance}{}.
\end{caveatbox}

Wave~6 brings the endpoint estate into the architecture's
authorisation surface: posture claims produced by managed devices
become an input to the policy decisions of \trchap{ch:policy}{},
unmanaged or BYOD devices are governed at the network attachment
layer of \trchap{ch:network}{}, and the institutional decommissioning
procedures are exercised against the new pipeline. The wave is also
the place where the architecture's most candid acknowledgement ---
the open-source gap on iOS and Android at scale (D11.3) ---
becomes operationally consequential.

%-----------------------------------------------------------------------------
\section{Goal and dependencies}
\label{sec:w6-goal}

The goal of Wave~6 is to deploy a unified posture schema across
the endpoint operating systems in production, to route the
resulting posture data into the policy plane established in Wave~4,
and to integrate the institutional decommissioning procedures with
the canonical event schema of \trchap{ch:observability}{}. The wave
does not require a single MDM product across all OSes; it requires
that the posture data produced by whatever products the institution
operates be normalised at the wave's boundary.

%-----------------------------------------------------------------------------
\section{Pre-requisites}
\label{sec:w6-prereq}

Beyond Waves~1 and 4, four operational conditions strengthen
Wave~6:

\begin{itemize}[leftmargin=1.5em,nosep]
  \item the endpoint estate is inventoried by operating system,
    by management state (managed, unmanaged, BYOD), and by user
    population (staff, student, contractor, visiting researcher);
  \item the existing endpoint-management products are documented
    against their licence terms and renewal trajectories;
  \item the boundary with the network layer's BYOD attachment
    policy is articulated, since posture verification of BYOD
    endpoints relies on network-side observation rather than on
    agent-based collection (D11.7);
  \item the unified posture schema (D11.6) is drafted and reviewed
    against the policy authoring cadence of Wave~4.
\end{itemize}

%-----------------------------------------------------------------------------
\section{Coexistence pattern}
\label{sec:w6-coexistence}

Wave~6 introduces posture signals into the policy plane in
\emph{advisory} mode before any enforcement decision depends on
them. The pattern parallels the shadow mode of Wave~4 but operates
on signal availability rather than on decision evaluation: the
posture claims are produced and routed to the policy layer, the
policy decisions are evaluated as if posture were enforced, but
no access is denied on posture grounds during the advisory
period.

\paragraph{Per-population transition.} Endpoint populations are
typically advanced through advisory-to-enforcing in cohorts:
managed staff devices first (operationally controlled), then
managed student devices, then BYOD with network-side enforcement.
Each cohort's transition is per-policy: a posture rule may be
enforcing for staff devices while still advisory for student
devices for an extended period.

\paragraph{Multi-product reconciliation.} For institutions with a
mixed MDM estate (a dedicated management product for the Apple
platforms, a separate suite for Windows and Android, and
configuration-management-plus-osquery tooling for Linux), the
principal engineering effort
of Wave~6 is the reconciliation of the products' native posture
formats into the unified schema. Each product's posture data is
mapped, the mapping is exercised in shadow, and the policy plane
ultimately consumes only the unified schema --- never the
product-specific signals.

%-----------------------------------------------------------------------------
\section{Trajectory: greenfield}
\label{sec:w6-greenfield}

For institutions in a greenfield position, Wave~6 deploys the
endpoint products and the unified schema together: there is no
incumbent MDM logic to reconcile, and the posture-to-policy
integration is built from the start. Indicative duration is six to
nine months for an institution of moderate scale.

%-----------------------------------------------------------------------------
\section{Trajectory: brownfield single-suite-centric}
\label{sec:w6-brownfield-ms}

For institutions whose endpoint management is dominated by a
single proprietary management suite, Wave~6 is the work of routing
that suite's compliance signals into the unified schema and of
reconciling its native authorisation with the policy plane.

\paragraph{Pattern.} The incumbent suite's compliance policies
remain in place for device-level enforcement (encryption, screen
lock, OS version, patch SLAs); their compliance signals are
exported through OpenTelemetry collectors deployed during Wave~2
and reach the policy plane in the unified schema. The conditional-access
expressions of the institution's incumbent identity suite that
depend on device compliance are exercised in shadow mode through
the equivalent OPA policies authored in Wave~4. Linux endpoints
--- typically present in environments characteristic of the
institution's research-computing profile even where a single
proprietary suite dominates --- are managed through
configuration-management and osquery tooling, with their posture
also normalised to the unified schema.

\paragraph{The iOS gap.} For institutions with a substantial iOS
estate, a proprietary suite's iOS coverage is the working choice;
there is no credible open-source alternative at scale at the time
of writing.
This is a known limitation of the architecture (D11.3 in
\trchap{ch:endpoint}{}). The
sovereignty consequence --- a Partial score on the technical and
institutional dimensions for the iOS portion of the estate ---
should be acknowledged in the wave-closure report rather than
elided.

\paragraph{Indicative duration.} Twelve to eighteen months,
dominated by per-population cohort transition.

%-----------------------------------------------------------------------------
\section{Trajectory: brownfield hybrid}
\label{sec:w6-brownfield-hybrid}

For institutions with a multi-OS estate combining Apple, Windows,
Linux and mobile platforms across organisational-unit boundaries, Wave~6 is
principally the work of \emph{normalisation}: producing a common
posture schema across products that were chosen at different times
and managed by different teams.

\paragraph{Pattern.} Each existing endpoint-management product
(a dedicated suite for macOS and iOS where present, a separate
suite for Windows and Android, configuration-management-plus-osquery
tooling for Linux, possibly a further unified-endpoint product in
some segments) continues to operate in its own domain. A
collector layer translates each product's posture format into the
unified schema; the policy plane consumes only the unified schema.
The institutional posture schema becomes the authoritative artefact;
the products become substitutable consumers of the schema.

\paragraph{What changes in Wave~6.} The posture data flow becomes
uniform across the institution; the policy plane no longer
branches on operating system for routine decisions (D11.6); the
decommissioning procedures become testable end-to-end against the
canonical event schema. The product mix itself is not changed by
this wave.

\paragraph{Indicative duration.} Twelve to eighteen months, with
the variability driven by the heterogeneity of the existing
product mix.

%-----------------------------------------------------------------------------
\section{Reversibility criteria}
\label{sec:w6-reversibility}

Wave~6 is reversible at the per-policy and per-population level
throughout the advisory-to-enforcing transition. Beyond that
transition, reversal is operationally possible but requires the
incumbent enforcement logic to be re-instated.

\begin{itemize}[leftmargin=1.5em,nosep]
  \item The advisory / enforcing toggle is per posture rule and per
    population; reversal of a posture rule is the deployment of a
    new policy bundle, not the redeployment of the endpoint
    products.
  \item Per-population rollback: if a population's transition to
    posture-enforced authorisation produces unworkable behaviour,
    that population is returned to advisory mode while the
    discrepancy is diagnosed.
  \item The unified posture schema can be revised; a schema change
    is treated as a governance event rather than an operational
    one (D12.2).
\end{itemize}

%-----------------------------------------------------------------------------
\section{Exit criteria}
\label{sec:w6-exit}

Wave~6 closes when the three Endpoint tests of the conformance
suite of \trchap{ch:conformance}{} (\S\ref{sec:cf-suite}) pass on
the unified posture pipeline, when the agreed populations are in
posture-enforced authorisation, and when the BYOD attachment
boundary with \trchap{ch:network}{} is verified.

\begin{itemize}[leftmargin=1.5em,nosep]
  \item Posture data from each managed OS arrives in the canonical
    schema; decommissioning of a representative endpoint executes
    end-to-end; MDM configuration profiles export to a documented
    format that an alternative MDM can consume on a representative
    subset. (Endpoint conformance subset.)
  \item Population coverage: the institutionally agreed populations
    are operating under posture-enforced authorisation; the
    remainder is in a documented schedule or exception register.
  \item BYOD boundary: the network-side BYOD posture verification
    of \trchap{ch:network}{} is consistent with the institution's
    BYOD acceptable-use policy.
\end{itemize}

\begin{realworldbox}[Endpoint management at large institutions]
\noindent Public documentation of endpoint architectures at
large institutions is uneven. Patterns reported in \loc{nren-name}'s community fora and at sector
conferences include mixed-OS estates combining a dedicated suite
for Apple platforms, a separate suite for Windows and Android, and
open tooling (configuration management, osquery, and macOS package
managers) for Linux and macOS. Specific institutional
configurations and the current
state of any cited deployment should be confirmed against the
institution's own published material before being cited as
precedent.
\end{realworldbox}

%-----------------------------------------------------------------------------
\section{Regulatory anchoring}
\label{sec:w6-regulatory}

\noindent Wave~6 brings managed endpoints under the same
attestation discipline as the rest of the architecture. It
contributes to \gls{gdpr}~Art.~32 (technical measures on the
device that processes personal data) and operationalises the
applicable national cybersecurity baseline regime's risk-management
measures --- the national transposition of the NIS2 Directive where
it applies --- specifically cyber-hygiene practices
(NIS2~Art.~21(2)(g)), asset management (Art.~21(2)(i)) and,
indirectly, multi-factor authentication (Art.~21(2)(j)) where the
endpoint carries the authentication factors. The wave operationalises 27001~A.8.1
(user endpoint devices), A.8.7 (protection against malware),
A.8.9 (configuration management) and contributes to NIST CSF
\textsf{PROTECT} (PR.PS --- platform security). Detailed mapping
in \trchap{ch:regulatory-mapping}{} \S\ref{sec:rm-gdpr},
\S\ref{sec:rm-nis2}, \S\ref{sec:rm-iso},
\S\ref{sec:rm-nist-csf}.
